\section{Related Work}
\label{sec:related}

\subsection{Invariant Wi-Fi Representations}

SHARP estimates a sparse channel impulse response and uses its strongest path as
a phase reference before constructing Doppler spectrograms
\cite{meneghello2023sharp}. Its representation is designed to transfer across
people and environments. Widar3.0 pursues the same general objective through a
physics-derived body-coordinate velocity profile, showing the value of removing
domain-dependent propagation effects before recognition
\cite{zheng2019widar}. Both systems therefore place substantial inductive bias
before the classifier rather than asking a generic classifier to discover
invariance from raw CSI.

\subsection{Phase-Aware CSI Processing}

Measured CSI contains amplitude variation and packet-dependent phase errors
caused by timing, carrier-frequency, sampling-frequency, and hardware offsets.
Such nuisance terms can dominate the much smaller phase changes induced by
motion. Ratnam \emph{et al.} formulate CSI preprocessing explicitly and analyze
how preprocessing choices affect sensing \cite{ratnam2024optimal}. The SHARP
dataset further documents the acquisition setting and its 80-MHz, 256-tone CSI,
of which 242 subcarriers are retained \cite{meneghello2023dataset}. These works
motivate treating phase correction as part of the learning problem, not as a
minor input normalization.

\subsection{Learning-Assisted Signal Processing}

Learning can reduce signal-processing cost without discarding known structure.
SLNet learns a spectrogram through a neural architecture tailored to wireless
sensing instead of regressing an unconstrained image
\cite{yang2023slnet}. LISTA unrolls sparse coding into a finite learned network,
preserving the optimization structure while learning efficient updates
\cite{gregor2010lista}. U-Net supplies the multiscale skip-connected backbone
used in our students \cite{ronneberger2015unet}, but it does not itself encode
complex phase equivariance or wireless propagation. The gap addressed here is
therefore narrower than generic end-to-end HAR: we test whether a convolutional
student can directly distill the output of a phase-aware preprocessing pipeline,
and compare that attempt with retaining a minimal explicit phase correction.
